\chapter*{Einleitung zum Gesamtwerk}
\addcontentsline{toc}{chapter}{Einleitung zum Gesamtwerk}

\begin{merksatz}
\boxtitel{Das Ziel}
Dieses Werk verfolgt ein einziges Ziel: einem erwachsenen Menschen, der die Mathematik verstehen will, den Weg von null bis zur Einführungsvorlesung zu ebnen --- ohne Lücken, ohne Sprünge und ohne die Formel, die alles verdirbt: \enquote{wie man leicht sieht}.

\begin{center}
\begin{minipage}{0.9\textwidth}
\textbf{Nichts ist selbstverständlich. Nichts ist trivial. Jeder Schritt wird erklärt.}
\end{minipage}
\end{center}

\schritt{Die Regel, ausbuchstabiert}
\begin{center}
{\small
\renewcommand{\arraystretch}{1.52}
\begin{tabular}{@{}P{2.8cm}P{5.8cm}@{}}
\toprule
\textbf{Keine Regel} & ohne Begründung. Warum minus mal minus plus ergibt, steht in Band 0 auf Seite drei.\\
\textbf{Keine Formel} & die vom Himmel fällt. Die $p$-$q$-Formel wird in vier Zeilen quadratischer Ergänzung entwickelt.\\
\textbf{Kein Beweis} & mit Lücke. \enquote{Analog} und \enquote{offensichtlich} kommen nicht vor.\\
\textbf{Keine Aufgabe} & ohne vollständigen Rechenweg --- samt Zahlenprobe am Ende.\\
\bottomrule
\end{tabular}}
\end{center}

\begin{center}
\begin{minipage}{0.88\textwidth}
\kursiv{Der Grund für diese Strenge ist nicht Pedanterie, sondern Erfahrung. \textbf{Wer eine Regel ohne Begründung lernt, hat nichts, woran er sich erinnern könnte.} Wer versteht, warum sie gilt, kann sie rekonstruieren, wenn er sie vergisst.}
\end{minipage}
\end{center}
\end{merksatz}

\section*{Wie das Werk entstanden ist}
\addcontentsline{toc}{section}{Wie das Werk entstanden ist}

\begin{merksatz}
\boxtitel{Drei Bauabschnitte}
Die Bände sind nicht in der Reihenfolge geschrieben worden, in der sie hier stehen.

\begin{center}
{\small
\renewcommand{\arraystretch}{1.55}
\begin{tabular}{@{}P{2.4cm}P{6.2cm}@{}}
\toprule
\textbf{Erster} & Ein Band über den Beweisbegriff; danach Zahlentheorie, lineare Algebra, Analysis, Stochastik, Numerik, Geometrie, Geschichte, Topologie, Vektoranalysis, Logik. Elf Bände --- und dann fehlte der wichtigste.\\
\textbf{Band 0} & Alle elf hatten stillschweigend vorausgesetzt, was in der Schule gelehrt und meist wieder vergessen wird. Dieser Band trägt daher die Nummer null: Er wurde an dieser Stelle geschrieben und gehört an den Anfang.\\
\textbf{Zweiter} & Algebra, Funktionentheorie, Maßtheorie mit Funktionalanalysis, Differentialgleichungen, Kombinatorik --- die Bände 11 bis 15. Damit war der Stoff der ersten vier Semester abgedeckt.\\
\textbf{Dritter} & Elf Bände: zehn, die die vorhandenen Gebiete anwenden --- Graphentheorie, Statistik, Biomathematik, Optimierung, Informationstheorie, Spieltheorie, maschinelles Lernen, Kausalität, Entscheidungstheorie, Ökonometrie --- und einer, der eine Grundlagenlücke schließt: die stochastischen Prozesse.\\
\textbf{Vierter} & Vier Nachtragsbände: zwei Anwendungen --- numerische Simulation (27) und dynamische Systeme (28) --- und zwei Grundlagen, die in der Lückenliste der ersten Fassung standen: algebraische Geometrie (29) und Darstellungstheorie (30).\\
\bottomrule
\end{tabular}}
\end{center}

\begin{center}
\begin{minipage}{0.88\textwidth}
\kursiv{Der dritte Abschnitt hat den Charakter des Werkes verändert, und Abschnitt 1.2 behandelt das ausführlich. \textbf{Aus einem Lehrgang wurde ein Lehrgang mit Anwendungsteil} --- und die Anwendungsbände greifen so stark ineinander, dass sie eine eigene Lesart verdienen.}
\end{minipage}
\end{center}
\end{merksatz}

\begin{merksatz}
\boxtitel{Warum ein Vorkurs nötig ist}
Die Abbruchquoten in der Mathematik und den mathematiknahen Fächern gehören zu den höchsten überhaupt; der weitaus größte Teil der Abbrüche geschieht in den ersten drei Semestern.

\begin{center}
\begin{minipage}{0.9\textwidth}
\textbf{Es fehlt selten am Verständnis für den neuen Stoff. Es fehlt an der Verfügbarkeit des alten.}
\end{minipage}
\end{center}

\begin{center}
\begin{minipage}{0.88\textwidth}
\kursiv{Wer beim Ableiten von $\frac{x^{2}-1}{x+3}$ noch über das Zusammenfassen von Brüchen nachdenken muss, hat keine Aufmerksamkeit mehr für die Quotientenregel übrig --- und schon gar nicht für den Gedanken dahinter. \textbf{Deshalb steht am Anfang dieses Werkes kein Kapitel über Mengen und Abbildungen}, wie es in Vorlesungsskripten üblich ist, sondern eines über Zahlen und das Rechnen mit ihnen.}
\end{minipage}
\end{center}
\end{merksatz}

\section*{Die drei Prinzipien}
\addcontentsline{toc}{section}{Die drei Prinzipien}

\begin{merksatz}
\boxtitel{Erstens: Herkunft statt Fertigware}
Jeder Begriff wird dort abgeholt, wo er entstanden ist.

\begin{center}
{\small
\renewcommand{\arraystretch}{1.5}
\begin{tabular}{@{}P{2.8cm}P{5.8cm}@{}}
\toprule
\textbf{Die Mengenlehre} & erscheint nicht als Fundament, sondern als das, was sie war: Cantors Werkzeug für ein technisches Problem der Fourieranalyse.\\
\textbf{Der Funktionsbegriff} & kommt nicht als Definition, sondern als Dirichlets Bruch mit Euler.\\
\textbf{Das Auswahlaxiom} & kommt mit Russells Socken und Zermelos Streit mit den französischen Analytikern.\\
\bottomrule
\end{tabular}}
\end{center}

\begin{center}
\begin{minipage}{0.88\textwidth}
\kursiv{Der Grund ist didaktisch, nicht antiquarisch. \textbf{Ein Begriff, dessen Problem man kennt, ist verständlich; ein Begriff ohne Problem ist Willkür.} Ein eigener Band --- Band 7 --- ist der Geschichte gewidmet und erzählt die Entwicklung im Zusammenhang, statt sie in Fußnoten zu verteilen.}
\end{minipage}
\end{center}
\end{merksatz}

\begin{warnung}
\boxtitelr{Zweitens: Die Fehler ernst nehmen}
In jedem Band gibt es Warnkästen. Sie enthalten keine Nebensächlichkeiten, sondern die Stellen, an denen erfahrungsgemäß etwas schiefgeht:
\[ -2^{2} \ne (-2)^{2}, \qquad (a+b)^{2} \ne a^{2}+b^{2}, \qquad \sqrt{a^{2}} = |a|, \qquad f^{-1} \ne \tfrac1f \]
\[ \frac{a}{b+c} \ne \frac ab + \frac ac, \qquad \log(a+b) \ne \log a + \log b, \qquad (fg)' \ne f'g' \]

\begin{center}
\begin{minipage}{0.88\textwidth}
\kursiv{Band 0 sammelt dreiunddreißig solcher Fehler in einem Katalog, jeweils mit falscher Fassung, richtiger Fassung und Kapitelverweis. \textbf{Fehler sind keine Zufälle. Sie folgen Mustern} --- fast alle beruhen darauf, dass man eine Operation für linear hält, die es nicht ist. Wer das Muster kennt, erkennt den Fehler wieder.}
\end{minipage}
\end{center}

\schritt{In den Anwendungsbänden ändert sich die Art der Fehler}
\begin{center}
\begin{minipage}{0.88\textwidth}
\kursiv{Dort geht es nicht mehr um Rechenfehler, sondern um Anwendungsfehler: ein Modell, das nicht passt; eine Voraussetzung, die niemand geprüft hat; ein Satz, der auf einen Fall angewandt wird, den er nicht abdeckt. \textbf{Die Warnkästen heißen dort \enquote{Der Bruchfall}} und stehen in jedem Kapitel.}
\end{minipage}
\end{center}
\end{warnung}

\begin{merksatz}
\boxtitel{Drittens: Bis an den heutigen Rand}
Jeder Band endet mit einem Kapitel \enquote{Aktueller Stand} --- was heute offen ist, woran gearbeitet wird, was sich in den letzten Jahren geändert hat.

\begin{center}
{\small
\renewcommand{\arraystretch}{1.46}
\begin{tabular}{@{}P{1.8cm}P{6.4cm}@{}}
\toprule
Band 1 & Was die Riemannsche Vermutung mit Primzahlen zu tun hat\\
Band 3 & Die Navier-Stokes-Gleichungen als Millenniumsproblem\\
Band 5 & Warum P gegen NP die Numerik betrifft\\
Band 8 & Die Poincaré-Vermutung und Perelmans Beweis\\
Band 10 & Woodins Programm zur Kontinuumshypothese\\
Band 17 & Die Replikationskrise und was aus ihr folgt\\
Band 19 & Wo lineare Optimierung an ihre Grenzen stößt\\
Band 20 & Kryptographie nach dem Quantenrechner\\
Band 21 & Warum Menschen nicht spielen, was die Theorie sagt\\
Band 22 & Warum tiefe Netze verallgemeinern, obwohl sie es nicht dürften\\
Band 23 & Warum dieselbe Tabelle mit entgegengesetzten Ursachen vereinbar ist\\
Band 24 & Warum ein Test nur zwischen zwei Schwellen etwas nützt\\
Band 25 & Warum beim fairen Spiel meist einer fast durchgehend führt\\
Band 26 & Warum zwei unabhängige Zeitreihen meist signifikant zusammenhängen\\
Band 27 & Warum ein Simulationsergebnis ohne Fehlerbalken keines ist\\
Band 28 & Warum man das Wetter nicht über zwei Wochen hinaus vorhersagen kann\\
Band 29 & Warum Gröbner-Basen entscheiden, welche Parameter ein Modell preisgibt\\
Band 30 & Warum die McKay-Vermutung nach fünfzig Jahren bewiesen ist\\
\bottomrule
\end{tabular}}
\end{center}

\begin{center}
\begin{minipage}{0.9\textwidth}
\textbf{Der Stand ist überall der August 2026.}
\end{minipage}
\end{center}

\begin{center}
\begin{minipage}{0.88\textwidth}
\kursiv{Nicht, weil man das für die Prüfung braucht, sondern weil Mathematik sonst wie ein abgeschlossenes Gebäude wirkt. \textbf{Sie ist es nicht.}}
\end{minipage}
\end{center}
\end{merksatz}

\section*{Was das Werk nicht ist}
\addcontentsline{toc}{section}{Was das Werk nicht ist}

\begin{warnung}
\boxtitelr{Vier ehrliche Einschränkungen}
\begin{center}
{\small
\renewcommand{\arraystretch}{1.55}
\begin{tabular}{@{}P{2.4cm}P{6.2cm}@{}}
\toprule
\textbf{Erstens} & Es ist kein Ersatz für das Üben. Wer die durchgerechneten Aufgaben nur liest, statt sie mitzurechnen, wird beim ersten eigenen Problem stehen. Papier und Stift sind Teil der Ausrüstung.\\
\textbf{Zweitens} & Es ist keine Vollständigkeitserklärung. Die zweite Werkhälfte und die Nachtragsbände haben viele Lücken geschlossen --- diskrete Mathematik, Statistik, Optimierung, algebraische Geometrie, Darstellungstheorie, numerische Simulation ---, andere bleiben offen: Lie-Gruppen und -Algebren systematisch, partielle Differentialgleichungen in der Tiefe, algebraische Zahlentheorie, Differentialgeometrie.\\
\textbf{Drittens} & Es ersetzt keine Vorlesung. Es soll sie verständlich machen.\\
\textbf{Viertens} & Es ist nicht neutral. Wo es mehrere Zugänge gibt, ist einer gewählt und begründet worden --- und die Begründung steht dabei.\\
\bottomrule
\end{tabular}}
\end{center}

\begin{center}
\begin{minipage}{0.88\textwidth}
\kursiv{Zur zweiten Einschränkung eine Bemerkung: \textbf{Die ursprüngliche Fassung dieser Liste nannte Algebra, Funktionentheorie, Maßtheorie, Differentialgleichungen und Kombinatorik als fehlend.} Sie sind inzwischen die Bände 11 bis 15; die zweite Fassung nannte algebraische Geometrie, Darstellungstheorie und numerische Simulation --- sie sind die Bände 29, 30 und 27. Eine Lückenliste ist ein Arbeitsplan, kein Bekenntnis.}
\end{minipage}
\end{center}
\end{warnung}

\begin{merksatz}
\boxtitel{Und ein Wort zur Zählung}
\begin{center}
{\small
\renewcommand{\arraystretch}{1.5}
\begin{tabular}{@{}P{2.4cm}P{2.4cm}P{2.6cm}@{}}
\toprule
\textbf{Bände} & \textbf{Anzahl} & \textbf{Seiten}\\
\midrule
$0$ bis $15$, $29$, $30$ und $B$ & $19$ & $3348$\\
$16$ bis $28$ & $13$ & $2696$\\
\midrule
\textbf{Zusammen} & $\mathbf{32}$ & $\mathbf{6044}$\\
\bottomrule
\end{tabular}}
\end{center}

\begin{center}
\begin{minipage}{0.88\textwidth}
\kursiv{Der vorliegende Rahmenband ist nicht mitgezählt. \textbf{Der Begleitband \enquote{Die Architektur der Gewissheit} trägt keine Nummer}, weil er als erster entstand und sich mit Band 10 teilweise überschneidet. Wer nur einen von beiden liest, nehme Band 10 für die Technik und den Begleitband für die Einordnung.}
\end{minipage}
\end{center}
\end{merksatz}
